\subsection{The threefold \texorpdfstring{$X_9$}{X9} and its transition}
\label{subsec:x9}

We recall the facts about $X_9$ used below; all of them are established
in~\cite{W26global,W26quantum}.

The resolution $\wideX_9$ is a smooth anticanonical hypersurface in a toric fourfold whose fan is
built on the polytope with vertices $v_0,\dots,v_8$ of~\cite[eq.~(3.1)]{W26quantum}. For a vertex $v_i$ we write $D_i$ for the
restriction to $\wideX_9$ of the corresponding toric divisor. Blowing down the exceptional
locus gives a morphism $\pi:\wideX_9\to X_9$ onto a Calabi--Yau threefold $X_9$ with three
singular points: two ordinary double points and one cone over the degree-seven del Pezzo surface
$S_7=\operatorname{Bl}_2\mathbb P^2$. The exceptional locus of $\pi$ consists of the two nodal
curves $C_1,C_2$ and a surface $E\cong S_7$, the divisor of an interior lattice point of a
two-dimensional face of the polytope. We write $e_1,e_2$ for the two exceptional curves of $E$
and $h$ for the pullback of the hyperplane class of $\mathbb P^2$, so that $h-e_1-e_2$ is the
third $(-1)$-curve of $E$.

The singular threefold $X_9$ has an explicit smoothing $X_{9,t}$~\cite[Sect.~3.1]{W26quantum}.
Near the three singular points, write the nodes as $xy-zw=u_i$, $i=1,2$, and let $s$ be the
coordinate on the reduced local deformation space of the del Pezzo cone
of~\cite[eqs.~(2.2), (2.3)]{W26quantum}, denoted $\beta$ there; in the $E_2$ theory it is the
Cartan component of the moment map of its $SU(2)$ flavour symmetry. The image of every global
first-order deformation of $X_9$ in these local parameters lies on the line
\begin{equation}
 u_1+s=u_2+s=0,
 \label{eq:x9-line}
\end{equation}
and the smoothing moves along it in the direction $(u_1,u_2,s)=t(-1,-1,1)$, smoothing all
three points at once~\cite[eqs.~(3.4), (3.9)]{W26quantum}; the statement about all global
first-order deformations follows from the proof of~\cite[Thm.~2.2]{W26quantum}, which applies to
the full space of local first-order deformations. The reason is that the two nodal curves and the
class $e_1-e_2$ of the del Pezzo surface satisfy $\iota_*(e_1-e_2)=C_1+C_2$ in
$H_2(\wideX_9;\ZZ)$, so the three local sectors are charged under the same two vector
multiplets. The transition changes the Hodge numbers from $(h^{1,1},h^{2,1})=(6,122)$ to
$(3,123)$~\cite[eqs.~(1.2), (3.5)]{W26quantum}: three vector multiplets are lost, two Higgsed by
the nodal and del Pezzo hypermultiplets and one because the Coulomb direction of the $E_2$ theory
disappears, and one neutral hypermultiplet appears.

\subsection{A face on which four curves contract}
\label{subsec:face}

Let $H_1,\dots,H_6$ be the integral basis of $H^2(\wideX_9;\ZZ)$ formed by the
generators of the ambient nef cone restricted to $\wideX_9$~\cite[eq.~(C.1)]{W26quantum},
and $\beta_1,\dots,\beta_6$ the dual curve classes. Writing $\iota:E\to\wideX_9$ for the inclusion, in
this basis
\begin{equation}
 C_2=\beta_1,\qquad C_1=\beta_5,\qquad \iota_*e_2=\beta_6,\qquad
 \iota_*e_1=\beta_1+\beta_5+\beta_6,\qquad \iota_*(h-e_1-e_2)=\beta_2,
 \label{eq:face-charges}
\end{equation}
and $H_3=D_1$, $H_4=D_0$ are the two classes pulled back from $X_9$; the divisors $D_0,D_1$
do not meet the exceptional locus. For $J$ in the K\"ahler cone the gauge kinetic matrix of the
vector multiplets is proportional to the Hodge metric $\int\omega_I\wedge*_J\omega_K$ on harmonic
two-forms. With $L(J)=(\kappa_{IKM}J^M)_{I,K=1}^{6}$ the full $6\times6$ matrix on
$H^2(\wideX_9)$, $v=LJ$ and $\operatorname{vol}=J^3/6$, it is
$\mathcal G=-L+vv^{\mathsf T}/(4\operatorname{vol})$~\cite[eqs.~(4.1), (4.6)]{W26quantum}. For curve classes $\gamma_1,\dots,\gamma_k$ we write $\operatorname{ann}(\gamma_1,\dots,\gamma_k)$
for the subspace of $H^2(\wideX_9;\mathbb R)$ on which they all vanish. Set
\begin{equation}
 F:=\{\,J=y_2H_2+y_3H_3+y_4H_4\ :\ y_2,y_3,y_4>0\,\}.
 \label{eq:face-def}
\end{equation}
Since $H_2\cdot e_i=0$ and $H_2\cdot(h-e_1-e_2)=1$, one has $H_2|_E=h$, and hence
$J|_E=y_2\,h$ for $J\in F$.

\begin{proposition}\label{prop:face}
$\overline F=\operatorname{Nef}(\wideX_9)\cap\operatorname{ann}(C_1,C_2,e_1,e_2)$, and
for $J\in F$:
\begin{enumerate}
\item an irreducible curve $C\subset\wideX_9$ has $J\cdot C=0$ if and only if
 $C\in\{C_1,C_2,e_1,e_2\}$; these four curves are pairwise disjoint and each has normal
 bundle $\mathcal O(-1)\oplus\mathcal O(-1)$;
\item every prime divisor has positive volume; in particular
 $\operatorname{vol}(E)=\tfrac12J^2\cdot E=y_2^2/2$;
\item the gauge kinetic matrix $\mathcal G$ is positive definite.
\end{enumerate}
\end{proposition}

\begin{proof}
Since $H_2,H_3,H_4$ are nef, $J\cdot C=0$ forces $H_2\cdot C=D_1\cdot C=D_0\cdot C=0$.
The classes $y_3D_1+y_4D_0$ with $y_3,y_4>0$ are exactly
$\pi^*\operatorname{Amp}(X_9)$. Indeed $\pi^*\operatorname{Nef}(X_9)
=\operatorname{Nef}(\wideX_9)\cap\pi^*H^2(X_9;\mathbb R)=\operatorname{cone}(D_0,D_1)$
\cite[Prop.~4.1]{W26quantum}, and every line bundle $L$ on $\wideX_9$ of degree zero on all curves in the
fibres of $\pi$ is pulled back from $X_9$: since $K_{\wideX_9}=0$ and every line bundle is $\pi$-big for
birational $\pi$, the relative basepoint-free theorem~\cite[Thm.~3.24]{KollarMori:1998}, with $\Delta=0$
and $a=1$, shows that $L^b$ and $L^{b+1}$ are $\pi$-globally generated for $b\gg0$; being of degree zero
on the connected fibres, they are trivial there and descend, as $\pi_*\mathcal O_{\wideX_9}=\mathcal O_{X_9}$
for the normal variety $X_9$, and hence so does $L=L^{b+1}\otimes L^{-b}$, as
in~\cite[proof of Thm.~3.7, Step~9]{KollarMori:1998}. By Kleiman's criterion~\cite{Kleiman:1966} the
descended class of $y_3D_1+y_4D_0$ is ample for $y_3,y_4>0$. Hence $C$ is contracted by $\pi$, that is
$C\subset C_1\cup C_2\cup E$. On $E$ one has $J|_E=y_2h$, and the irreducible curves of
$E$ with $h\cdot C=0$ are $e_1$ and $e_2$. The two nodes are distinct points of $X_9$ away
from the image of $E$, and $e_1\cdot e_2=0$ on $E$, so the four curves are disjoint. For
$e_i$ one has $N_{e_i/E}=\mathcal O(-1)$ and $E\cdot e_i=K_E\cdot e_i=-1$, and the
extension of normal bundles splits; the $C_i$ are the exceptional curves of small
resolutions of ordinary double points. With~\eqref{eq:face-charges} this proves (1) and
$\overline F\subseteq\operatorname{Nef}(\wideX_9)\cap\operatorname{ann}(C_1,C_2,e_1,e_2)$;
the reverse inclusion holds because $\beta_2$, $\beta_3$ and $\beta_4+\beta_5+\beta_6$ are
effective (their genus-zero Gopakumar--Vafa invariants are $1$, $-2$ and $10$), which forces
$y_2,y_3,y_4\ge0$.

For (2), a class $J\in F$ with rational coordinates is nef and big, so by the
basepoint-free theorem~\cite[Thm.~3.3]{KollarMori:1998} some multiple defines a
birational morphism $\pi_F:\wideX_9\to X_F$ with $J=\pi_F^*A$, $A$ ample. If a prime divisor
$D$ had $J^2\cdot D=A^2\cdot\pi_{F*}D=0$, then $\dim\pi_F(D)\le1$ and $D$ would lie in the
exceptional locus of $\pi_F$, which by (1) is the union of four curves. At an irrational point
write $J=\sum_ia_iJ_i$ with $a_i>0$ and rational $J_i\in F$; then
$J^2\cdot D=\sum_{i,j}a_ia_jJ_i\cdot J_j\cdot D\ge\sum_ia_i^2J_i^2\cdot D>0$, because the mixed
terms are nonnegative for nef $J_i,J_j$ and effective $D$. The value of $\operatorname{vol}(E)$
follows from $J|_E=y_2h$.

For (3), since $v=LJ$ we have $v^{\mathsf T}L^{-1}v=J^{\mathsf T}LJ=6\operatorname{vol}$, and the
matrix determinant lemma gives
$\det\mathcal G=\det(-L)\bigl(1-v^{\mathsf T}L^{-1}v/(4\operatorname{vol})\bigr)=-\tfrac12\det L$, where
$\det(-L)=\det L$ because the size is even. In the coordinates $(y_2,y_3,y_4)$,
\begin{equation}
 \det L(J)=-y_2(2y_3+3y_4)(y_2+2y_4)\,P(y_2,y_3,y_4),
 \label{eq:detL}
\end{equation}
with $P$ the polynomial with positive coefficients printed in~\eqref{eq:app-detL}. So
$\det\mathcal G>0$ on $F$. On $\overline F\smallsetminus\mathbb R_{\ge0}H_2$ the volume is positive
and $\mathcal G$ is a limit of Hodge norms from the K\"ahler cone, hence positive semidefinite, and it
is positive definite wherever its determinant is nonzero.
\end{proof}

Four flop walls meet along $F$, one for each of the curves in (1). When the four curves are
used as the curves $C_1,\dots,C_4$ of Section~\ref{subsec:setting} we put $C_3=e_1$ and
$C_4=e_2$.

\subsection{The contraction and its smoothing}
\label{subsec:face-smoothing}

The contraction $\pi_F:\wideX_9\to X_F$ of the proof of Proposition~\ref{prop:face} has
four ordinary double points, the images of $C_1,C_2,e_1,e_2$, and the classes of these
curves satisfy the single relation
\begin{equation}
 [C_1]+[C_2]-[e_1]+[e_2]=0 ,
 \label{eq:face-relation}
\end{equation}
all of whose coefficients have absolute value one. By Friedman's
criterion~\cite{Friedman:1986} and the unobstructedness of its deformations~\cite{Gross:1997},
$X_F$ is smoothable, with a one-dimensional space of local
smoothing directions compatible with~\eqref{eq:face-relation}.

Write $u_{C_i},u_{e_i}$ for the local smoothing parameters of the four double points of $X_F$
(for a node $xy-zw=u$), and $(u_1,u_2,s)$ for the local parameters of $X_9$ of
Section~\ref{subsec:x9}.

\begin{proposition}\label{prop:face-smoothing}
A smoothing of $X_F$ in the direction permitted by~\eqref{eq:face-relation} has general fibre a
smoothing of $X_9$ in the direction $(u_1,u_2,s)=t(-1,-1,1)$
of Section~\ref{subsec:x9}; in particular it is deformation equivalent to $X_{9,t}$.
\end{proposition}

\begin{proof}
The morphism $X_F\to X_9$ contracts the image $\overline E\cong\mathbb P^2$ of $E$ to a
rational singularity, so deformations of $X_F$ blow down to deformations of
$X_9$~\cite[Prop.~11.4]{KollarMori:1992}. The permitted direction is
$(u_{C_1},u_{C_2},u_{e_1},u_{e_2})=t(-1,-1,1,-1)$. The two nodes of $X_9$ are the same germs as
the nodes $C_1,C_2$ of $X_F$, so the induced deformation has $u_i=u_{C_i}=-t\neq0$. The local
image of every global first-order deformation of $X_9$ lies on the line~\eqref{eq:x9-line}, by
the proof of~\cite[Thm.~2.2]{W26quantum}, so $s=t+O(t^2)\neq0$ for small $t\neq0$ as well. Hence
the induced deformation smooths both nodes and the del Pezzo cone point, and its general fibre
$X_{9,t'}$ is smooth. The induced morphism $f:X_{F,t}\to X_{9,t'}$ between smooth Calabi--Yau
threefolds satisfies $K_{X_{F,t}}=f^*K_{X_{9,t'}}$, so it is crepant; a crepant exceptional divisor over a smooth
target is impossible, so $f$ is small; and a small birational morphism onto a smooth, hence
$\mathbb Q$-factorial, variety is an isomorphism.
\end{proof}

By Proposition~\ref{prop:face}(1), and the vanishing of multiple-cover and higher-genus
contributions of isolated $(-1,-1)$-curves (Appendix~\ref{app:data}), the charged massless
spectrum at a point of $F$ consists
of four hypermultiplets, one for each curve, whose charge matrix has rank three and Smith
invariants $(1,1,1)$. Condensing them Higgses three vector multiplets and leaves one neutral
hypermultiplet, which matches $(h^{1,1},h^{2,1})=(6,122)\to(3,123)$.
